% 第6章: 3フェーズモデル
\section{3フェーズモデル}

%% ── Q1: loop()の整理 ──────────────────────
\begin{qabox}{loop()の中身はどう整理するの？}

ここまでで、モジュール・インターフェース・SystemDataの3つの道具が揃いました。最後のピースは「loop()の実行順序」です。

loop()の中身を3つのフェーズに分けます。これを「3フェーズモデル」と呼びます。

\begin{center}
\begin{tikzpicture}[
  phase/.style={rectangle, draw=#1, fill=#1!8, minimum width=50mm,
    minimum height=12mm, font=\bfseries, rounded corners=2mm},
  arr/.style={-{Stealth[length=4mm]}, very thick, headerblue},
]
  \node[phase=accentblue] (input) {1. 入力フェーズ};
  \node[phase=noteyellow, below=10mm of input] (logic) {2. ロジックフェーズ};
  \node[phase=okgreen, below=10mm of logic] (output) {3. 出力フェーズ};

  \draw[arr] (input) -- node[right, font=\small]{センサー値を読む} (logic);
  \draw[arr] (logic) -- node[right, font=\small]{判断・加工する} (output);
\end{tikzpicture}
\end{center}

工場の流れ作業に例えると分かりやすいです。

\begin{enumerate}
  \item \textbf{入力フェーズ}（検品係） — センサーからデータを読み取り、SystemDataに書き込む
  \item \textbf{ロジックフェーズ}（判断係） — SystemDataを見て、どう動くか判断する
  \item \textbf{出力フェーズ}（作業係） — SystemDataの指示に従い、LEDやモーターを制御する
\end{enumerate}

\begin{goodexample}{3フェーズのloop()}
\begin{lstlisting}
// モジュール配列
IModule* inputModules[] = { &tempSensor };
IModule* outputModules[] = { &led };
const int INPUT_COUNT = 1;
const int OUTPUT_COUNT = 1;

void loop() {
    // 1. 入力フェーズ — センサー値をSystemDataに書き込む
    for (int i = 0; i < INPUT_COUNT; i++) {
        if (inputModules[i]->enabled) {
            inputModules[i]->updateInput(systemData);
        }
    }

    // 2. ロジックフェーズ — データを見て判断する
    applyPattern(systemData);

    // 3. 出力フェーズ — 判断結果に従い制御する
    for (int i = 0; i < OUTPUT_COUNT; i++) {
        if (outputModules[i]->enabled) {
            outputModules[i]->updateOutput(systemData);
        }
    }
}
\end{lstlisting}
\end{goodexample}

\end{qabox}

%% ── Q2: 順番が大事な理由 ──────────────────
\begin{qabox}{なんで順番が大事なの？}

「入力→ロジック→出力」の順番を守らないと、思わぬバグが発生します。

\begin{badexample}{出力が先に実行される場合}
\begin{lstlisting}
void loop() {
    // 出力が先！
    led.updateOutput(systemData);   // 前回のデータで光ってしまう

    // 入力が後
    tempSensor.updateInput(systemData);  // 今回のデータを読む

    applyPattern(systemData);  // 今回のデータで判断するが...
    // → LEDにはもう反映されない（次のloopまで待つ）
}
\end{lstlisting}
\end{badexample}

出力が先だと、LEDは常に「1ループ前のデータ」で動作します。温度が急上昇しても、LEDの反応が1ループ分遅れます。通常は気にならない程度ですが、リアルタイム性が求められるシステムでは致命的です。

また、ロジックフェーズをスキップして入力モジュールが出力モジュールのデータを直接書き換えると、「なぜこのLEDが光っているのか」の理由がコード上で追えなくなります。

\begin{cautionbox}[クロス参照の禁止]
入力モジュールが出力モジュールのデータを読み書きしたり、その逆をしたりすることを「クロス参照」と呼びます。これは禁止です。

\begin{itemize}
  \item 入力モジュール → 自分のDataに書き込むだけ
  \item 出力モジュール → 自分のDataを読み取るだけ
  \item データの受け渡し → ロジックフェーズが担当
\end{itemize}
\end{cautionbox}

\end{qabox}

%% ── Q3: ロジックフェーズ ──────────────────
\begin{qabox}{ロジックフェーズでは何を書くの？}

ロジックフェーズは「入力データを見て、出力データを決める」場所です。

\begin{goodexample}{applyPattern()の例}
\begin{lstlisting}
void applyPattern(SystemData& data) {
    // 温度が30度を超えたらLEDを点灯
    if (data.temp.temperature > 30.0) {
        data.led.state = true;
    } else {
        data.led.state = false;
    }
}
\end{lstlisting}
\end{goodexample}

ポイントは以下の通りです。

\begin{itemize}
  \item ハードウェアを直接操作しない（\texttt{digitalWrite}等を呼ばない）
  \item SystemDataの値を読んで、SystemDataの値を書き換えるだけ
  \item 「なぜLEDが光るのか」のルールがこの関数に集約される
\end{itemize}

ロジックが複雑になっても、この関数を読めば動作の全体像が分かります。

\begin{notebox}[第1章のコードがこう変わった]
第2章では約40行のベタ書きでしたが、3フェーズモデルで整理すると次のような構成になります。

\begin{itemize}
  \item \texttt{LedModule} — LEDの制御（約20行）
  \item \texttt{TempSensorModule} — 温度の読み取り（約20行）
  \item \texttt{applyPattern()} — ルール（約5行）
  \item \texttt{main.cpp} — 配列定義 + 3フェーズループ（約30行）
\end{itemize}

合計行数は増えますが、各ファイルが短く、役割が明確です。
\end{notebox}

\end{qabox}
